% Generated from project outputs. Do not edit by hand.
\begin{table}[!htbp]
\centering
\begin{threeparttable}
\caption{Capacity-curve sensitivity: age 2016-2019 retuned}
\label{tab:robustness-age-2016-2019-retuned}
\small
\setlength{\tabcolsep}{4pt}
\renewcommand{\arraystretch}{1.15}
\begin{tabularx}{\linewidth}{@{}Xrrrrr@{}}
\toprule
Variant & Span & Q500 & Q1500 & Q3000 & Q5000 \\
\midrule
Main (14-60) & 0.300 & 667.5 & 1,774.5 & 3,148.0 & 5,403.0 \\
v1 (13-59) & 3.000 & 473.9 & 1,685.9 & 3,266.5 & 4,930.6 \\
v2 (15-61) & 0.275 & 688.0 & 1,780.0 & 3,153.3 & 5,474.0 \\
v3 (12-58) & 3.000 & 480.5 & 1,688.5 & 3,268.8 & 4,943.3 \\
v4 (16-62) & 0.225 & 774.9 & 1,776.3 & 3,143.7 & 5,695.1 \\
\bottomrule
\end{tabularx}
\begin{tablenotes}[flushleft]
\footnotesize
\item Entries are predicted income at t+1 (quetzales); column headings give income at t. General sample, labor income, population-weighted LOESS. -- indicates no supported prediction. Retuned selects a span for each cutoff; fixed\_baseline holds the main-sample span fixed.
\end{tablenotes}
\end{threeparttable}
\end{table}
